\subsection*{先独立规定一次购买应该怎样回答}

在检查 SQL 的执行步骤以前，先给一个只按购买含义工作的顺序对象。顺序对象每次完整处理一项动作，前一项结束后才处理下一项；它为并发历史提供可比较的业务规则。本章仍假定会话资格有效、记录不回收、报价在本地可验证，账户只有购买支出。

对象状态由当前余额 \(b\) 和终态映射 \(M\) 组成。映射按完整 \(S\) 保存原规范化版本、输入、购买号和结果。报价函数 \(q(X)\) 从规范化输入 \(X\) 中的商品、数量、币种及报价号取得应付金额 \(a\)；例如 Q7 对应 \code{paper-001}、数量 1、CNY 3000。这个金额在检查时从输入与报价规则独立计算，不能从实现已经写出的扣款倒推。

对尚未出现在 \(M\) 中的新 \(S\)，对象先验证首次成交依据。若 \(b\geq a\)，它把余额变为 \(b-a\)，建立属于该 \(S\) 的 P、O、D，保存相同商品、币种、金额和当次余额；若 \(b<a\)，余额保持原样，保存所需金额和余额不足原因，不建立 O、D。从 10000 购买 Q7，合法结果是扣 3000、剩 7000；从 2000 购买同一报价，合法结果是余额不足。10000 时稳定地拒绝，或者扣 4000 且流水也写 4000，都违反这项业务规则，即使它们没有破坏身份唯一或资金守恒。

对已有 \(S\)，对象按保存的规范化版本比较输入。同参返回 \(M(S)\) 的原结果，余额、关联和终态全部保持；异参返回冲突，也保持这些状态。因此，旧 Q7 后来到期不会使原成功重新判价，同参也不能被任意报告成冲突。查询是单独的只读动作：在其观察点，有终态就返回它，没有就返回未找到，状态保持不变。

\begin{figure}[htbp]
\centering
\begin{tikzpicture}
\node[booknode,text width=3.1cm] (input) at (1.7,2.1) {独立输入规则\\Q7、数量 1、CNY\\应付金额 3000};
\node[booknode,text width=3.1cm] (state) at (6,2.1) {顺序对象\\余额 b 与终态映射 M\\按完整 S 处理};
\node[bookresult,text width=3.1cm] (expected) at (10.3,2.1) {独立预期\\新 S 判断资金\\旧 S 返回原结果};
\draw[bookarrow] (input)--(state);
\draw[bookarrow] (state)--(expected);
\node[booknode,text width=4.6cm] (impl) at (3.5,-.25) {被检查的实现迁移\\身份、余额、效果、原结果};
\node[booknode,text width=4.6cm] (compare) at (9,-.25) {逐项比较投影和输出\\核对金额、身份、输入、结论};
\draw[bookarrow] (impl)--(compare);
\draw[bookalert] (expected.south)--(compare.north);
\node[booknote,text width=11cm] at (6,-1.65) {首次提交按原尝试的输入求预期；查询在快照点读取\\预期不取自被检查代码的成功判断，也不从实际流水反推金额};
\end{tikzpicture}
\caption{独立顺序规格先给出合法业务决定，再核对每个实现迁移与返回内容}
\label{fig:ch06-independent-spec}
\end{figure}


把并发执行接到这个对象，需要说明每项动作放在什么位置。完整逻辑购买的区间从客户端保存 S 并开始提交，延续到最终获知终态，跨越 R1 超时和 R2 恢复；通信尝试结束不关闭这个区间。新成功与余额不足在本章条件下取完整提交点，重放与异参比较取读取旧事实的位置，查询取权威 SELECT 的快照开始。每个位置都落在相应调用区间内；一项已经完成后才开始的另一项必须排在其后。未回应但已经提交的购买可以补入顺序历史，提交前回滚的尝试不增加购买效果。

例如 R1 发出后超时，数据库随后提交 S1 到 7000，R2 才返回原成功，这是一项跨通信尝试的购买。若 Q1 的快照早于该提交，它可以稍后返回 404；若成功已经返回后 Q2 才开始，Q2 必须读到仍保留的 S1 终态。前一种查空可以放在提交前，后一种空结果不能满足这个实际先后关系。

\subsection*{再写清楚要一直保持的关系}

同一 \(S\) 至多扣一次，仍不足以描述一个正确购买。例如服务可以只扣一次 4000，或者只扣一次 3000 却把结果指向别人的订单。要判断每次中断后的状态是否允许，需要把数量、归属、金额与返回内容一起表达。把这类在每个可达的已提交状态中都必须成立的条件称为不变量；可达表示从规定初始状态出发，按允许步骤实际能够走到。

先用数值说明资金条件。从初始 10000 出发，S1 成功扣 3000，S2 成功扣 3000，账户应为 4000；S1 再恢复，流水没有增加，余额保持 4000。为了同样检查更多购买，记初始余额为 \(b_0\)，当前已提交余额为 \(b\)，已提交扣款流水的集合为 \(D\)。对其中一条流水 \(d\)，用 \(a(d)\) 表示正整数金额。只包含购买扣款的模型要求：
\[
 b=b_0-\sum_{d\in D}a(d),\qquad b\geq0.
\]
\(\sum\) 表示把集合内每条流水的金额相加。两条各 3000 的 \(D\) 得到总额 6000，代回便是 \(b=10000-6000=4000\)。这个等式把余额与全部实际流水相联；若没有流水却私自改了余额，或者留下流水却没有扣余额，都无法通过检查。真实系统如果加入充值，须为每笔充值保存可核对的记录，并在右边增加这些记录的正金额项，再检查扩展后的守恒关系。

接着表达身份与业务数量。沿用第四、五章的 \(P(S)\) 与 \(N_T(S)=|P(S)|\)，身份条件是 \(N_T(S)\leq1\)，且每一行都已具有完整终态。属于 S 的扣款条数为 \(N_D(S)\)：成功终态恰好关联一张 O 与一条 D，因而为 1；拒绝终态没有 O、D，因而为 0。O、D 的 \code{purchase\_id} 都为对应 P，D 的 \code{order\_id} 为这张 O，三者币种和成交金额相同。拆单和分期要另定义关系。

最后检查结果。\(R(S)\) 表示已提交记录保存的业务结果；\(\pi_B(r)\) 表示从一次响应 \(r\) 中取出业务字段的投影。投影的作用是明确比较范围：保留终态、原编号、商品、数量、报价、币种、金额、当次余额与结果版本，略去本次请求编号和发送时间。对于权限与输入均符合要求的同参恢复，要求 \(\pi_B(r)=R(S)\)。两次恢复的 HTTP 状态可以为 200，首次新成功可以为 201；每次响应使用本次通信的包装字段，业务字段则持续等于原来的 \(R(S)\)。

\begin{figure}[htbp]
\centering
\begin{tikzpicture}
\node[booknode,text width=2.4cm] (s) at (1.4,2.3) {完整 S\\至多一份 P};
\node[bookresult,text width=2.6cm] (p) at (5.1,2.3) {成功 P1\\输入、币种、3000};
\node[booknode,text width=2.6cm] (o) at (9.2,3.1) {恰好一张 O1\\归属 P1};
\node[booknode,text width=2.6cm] (d) at (9.2,1.5) {恰好一条 D1\\归属 P1、O1};
\draw[bookarrow] (s) -- (p);
\draw[bookarrow] (p) -- (o);
\draw[bookarrow] (p) -- (d);
\node[booknode,text width=4.5cm] (r) at (2.6,-.2) {保存 \(R(S)\)：实际 P1、O1、D1\\每次同参恢复 \(\pi_B(r)=R(S)\)};
\node[booknode,text width=4.5cm] (money) at (8.4,-.2) {\(b=b_0-\sum_{d\in D}a(d)\)\\余额非负，扣款均有终态};
\draw[bookarrow] (p.south) -- (r.north);
\draw[bookarrow] (d.south) -- (money.north);
\end{tikzpicture}
\caption{数量、资金、归属与结果投影共同约束成功记录}
\label{fig:ch06-invariants}
\end{figure}


这些条件还需双向核对。成功终态引用的效果确实存在；反过来，每条购买扣款都能沿其购买回到唯一完整 \(S\) 与原结果。只查前一个方向，会放过没有完成记录的孤立扣款。拒绝分支则保留输入与原因，供恢复证明它曾明确结束；只保存 \code{handled=true}，无法满足原结论可恢复的要求。

\subsection*{把前提逐项接到允许步骤}

证明需要一个确定的协议范围。调用方为同一确认动作使用同一 \(S\)，为不同确认动作使用不同 \(S\)，且恢复保留原输入。若调用方把一次购买拆成 K1、K2，服务看到的是两项合法身份，唯一约束无法替它重建丢失的意图。如果把两次确认都写成 K1，服务合并恢复也符合自身规则，错误已经发生在身份分配处。

数据库方面，四列非空组合唯一约束可靠执行，所有业务写入经过这项仲裁，且只有实际插入身份的事务可以首次扣款。账户条件更新依照 PostgreSQL 18 Read Committed 的等待与重新判断规则执行；身份、余额、订单、流水与终态处于同一原子事务。原子表示外部已提交状态共同出现或共同撤销，事务的中间写入不单独发布。已确认提交在故障后继续保留，恢复不会切到遗失 P1 的旧历史。

写入与读取也承担前提。代码验证 Q7 的本地成交依据，以同一个 3000 和同一币种写入全部关系，从实际生成的编号构造结果。终态后身份、输入与原结果不被任意改写。恢复按完整 \(S\) 读取权威数据，先读原输入版本再比较，使用保存的结果版本输出。唯一约束限制记录条数，外键限制引用存在与归属，跨表金额相等仍需要正确写入路径维持。

Lee 等人的 RIFL 围绕可重试远程调用建立身份、完成记录和结果查找\cite{rifl}。调用方失联以后，重发须接回原执行。RIFL 要求完成记录与业务变化原子地持久保存，重试还必须找到同一份记录。本章的共同事务与完整 \(S\) 读取分别落实这两项关系。论文的迁移、租约与回收设计范围更广，当前证明采用已经具体实现的本地条件。

\subsection*{从基础状态到每一次迁移}

证明不变量可以按执行步数作归纳：先证明零步时成立，再证明只要当前状态成立，任一允许的下一步都使它继续成立。这样就覆盖任意有限步的允许执行，而无需逐条列出所有故障次数。这里的归纳对象是协议抽象步骤；在已声明的数据库唯一性、原子性与持久性前提下，逐项检查购买协议允许的迁移。

基础状态是 \(b=b_0=10000\)，\(D\)、订单和购买集合均为空。金额总和为零，余额等式成立；每个 \(S\) 的购买数量为零，成功与拒绝的关系没有需要核对的对象，结果条件也没有已发布响应需要满足。权限与格式在准入前检查，输入按适用版本整理；首次报价在取得 S 后、资金与成交写入前验证，已有 S 则先恢复原绑定。

身份插入、余额修改、订单插入、流水插入和终态填写，都先只修改某个事务的内部状态。对于其他观察者所见的已提交业务状态，这些步骤尚未发布变化，因此已提交不变量继续成立。内部状态当然可以暂时缺订单或缺结果；协议要求在它成为可见事实以前补齐或回滚。把已提交不变量机械施加到内部第一步，会把合法的组内构造过程也误判成错误。

回滚或提交前崩溃撤销内部组，不增加已提交购买与流水，\(b\) 也恢复到原有值。其他事务此前若被锁阻挡，之后可以继续，但必须重新走同一仲裁与条件更新。无论一项尝试在内部执行过多少步，只要没有发布提交，就没有为已提交集合增加半笔购买。

成功提交是第一个需要检查全部关系的步骤。组合唯一约束保证原先不存在另一条已提交同 \(S\) 购买，也禁止两个竞争者都成功提交同一 \(S\)。取得者写入一份 P、一份 O 与一份 D，并从实际编号构造原结果。条件更新将旧余额 \(b\) 减去已验证金额 \(a\)，同时 \(D\) 新增金额 \(a\)；资金等式两边同时减少 \(a\)。更新条件要求旧 \(b\geq a\)，所以新余额 \(b-a\geq0\)。所有新增关联与终态共同发布，成功所需关系成立。

\begin{figure}[htbp]
\centering
\begin{tikzpicture}
\node[booknode,text width=9.5cm] (base) at (5.5,4.2) {基础：余额等于初始值，购买、订单、流水为空；不变量成立};
\node[booknode,text width=9.5cm] (current) at (5.5,2.9) {归纳假设：某个当前已提交状态满足全部不变量};
\draw[bookarrow] (base) -- (current);
\node[booknode,text width=3cm] (internal) at (1.6,1.2) {内部步骤或回滚\\外部状态保持};
\node[bookresult,text width=3cm] (success) at (5.5,1.2) {成功或拒绝提交\\逐项核对新增关系};
\node[booknode,text width=3cm] (read) at (9.4,1.2) {重放、冲突、查询\\保留原业务状态};
\draw[bookarrow] (current) -- (internal);
\draw[bookarrow] (current) -- (success);
\draw[bookarrow] (current) -- (read);
\node[booknode,text width=9.5cm] (next) at (5.5,-.6) {各类允许迁移后的下一状态仍满足不变量；继续应用归纳步骤};
\draw[bookarrow] (internal) -- (next);
\draw[bookarrow] (success) -- (next);
\draw[bookarrow] (read) -- (next);
\end{tikzpicture}
\caption{从基础建立归纳，再对任意当前状态逐类检查下一步}
\label{fig:ch06-induction}
\end{figure}

\begin{figure}[htbp]
\centering
\begin{tikzpicture}
\node[booknode,text width=4.4cm] (before) at (2.5,1.5) {提交前\\\(b=7000,\ \sum a(d)=3000\)\\\(7000=10000-3000\)};
\node[bookresult,text width=4.4cm] (after) at (8.7,1.5) {提交后\\\(b'=4000,\ \sum a(d')=6000\)\\\(4000=10000-6000\)};
\draw[bookarrow] (before) -- (after);
\node[booknote] at (5.6,2.7) {同一事务扣 3000};
\node[booknote,text width=10.5cm] at (5.6,-.2) {余额减去 \(a\)，流水金额总和增加 \(a\)；\\旧余额至少为 \(a\)，所以新余额仍非负};
\end{tikzpicture}
\caption{成功提交同时改变余额与流水，保持同一守恒等式}
\label{fig:ch06-conservation-step}
\end{figure}


拒绝提交增加一条带输入与原因的 P，没有增加 O、D，也没有改变 \(b\)，资金等式保持。身份约束仍限制同 \(S\) 最多一条，原结果说明所需 3000 与余额不足，不声称有一笔扣款。此前约定只有购买支出，余额在条件判断与拒绝提交之间保持或下降，所以在选定的提交点仍可解释这项拒绝。

同参恢复在身份冲突后读原终态，不执行余额分支，不增加 P、O、D，资金与数量保持。按原版本解释结果使 \(\pi_B(r)=R(S)\)。异参冲突只结束当前尝试，同样不改动原记录。只读查询也不改变业务状态；其返回事实受语句快照限定，查空以后继续受约束的身份路径。

这些步骤穷尽当前协议对外发布业务变化的方式。因而从基础状态出发，任意次数的内部执行、回滚、竞争、提交与恢复，都保持列出的已提交不变量。若新增后台直接扣款、终态覆盖或记录删除，允许迁移集合便发生变化，须为新步骤补上保持论证。

\subsection*{从记录上限推出不重复生效}

现在可以把同 \(S\) 的成功效果与身份数量接起来。每份成功扣款必须关联一条购买，每条购买都关联完整 \(S\)；一条购买最多一条扣款，且同一 \(S\) 最多一条购买。因此同一 \(S\) 的已提交成功扣款数量至多为一。重复尝试即使在服务入口执行许多次，也无法在保持这些关系的同时增加第二份效果。

每个条件都有反例说明其必要作用。去掉 \(S\) 唯一性，可以有 P1 与 P2 各自合法地关联一条流水；去掉一笔购买一条扣款的规则，可以在 P1 下插入 D1 与 D2；允许扣款绕过购买，则即使 P1 唯一，也能产生第二条找不到 \(S\) 的支出；分开提交效果与结果，崩溃可留下没有身份证据的旧效果。这些错误都无法只靠限制余额非负发现，10000 足够重复扣两次 3000。

同 \(S\) 的余额不足也受同一身份上限约束。它已经提交以后，重试读回原拒绝，后来的资金条件变化不使这一项操作再进入首次判断。成功与拒绝在这里都是完成结论，只有成功分支形成扣款。至多一份终态和至多一份成功效果需要分别计数，拒绝的资金支出为零。

\subsection*{不错误合并与原结果一致分别依赖什么}

不同 \(S\) 的分离可以从两次同商品购买推导。S1 与 S2 的唯一键不同，规范化输入相同不会引起身份冲突。从 10000 开始，两笔各 3000 都可以通过条件更新，最终留下两个购买、两个订单与两条流水，余额 4000。它们在账户行上排队，影响先后与当次余额，却没有被身份查询认成同一条操作。

这里能够证明的是协议没有用商品或参数相等主动合并两项身份。若初始余额为 5000，第二项可能余额不足；若数据库永久不可用，它也可能长期未完成。不同身份各自有资格进入判断，成功仍取决于各自执行时的业务条件。要承诺它们最终各有终态，还需要下一节的进展条件。

\begin{figure}[htbp]
\centering
\begin{tikzpicture}
\node[booknote] at (5.6,3.5) {两次独立确认：S1、S2，商品与 Q7 相同，各需 3000};
\node[booknode,text width=2.8cm] (a) at (1.6,2.3) {初始 10000};
\node[bookresult,text width=2.8cm] (b) at (5.6,2.3) {S1 成功\\余额 7000};
\node[bookresult,text width=2.8cm] (c) at (9.6,2.3) {S2 成功\\余额 4000};
\draw[bookarrow] (a) -- (b);
\draw[bookarrow] (b) -- (c);
\node[booknode,text width=2.8cm] (d) at (1.6,.3) {初始 5000};
\node[bookresult,text width=2.8cm] (e) at (5.6,.3) {S1 成功\\余额 2000};
\node[booknode,text width=2.8cm] (f) at (9.6,.3) {S2 拒绝\\余额仍 2000};
\draw[bookarrow] (d) -- (e);
\draw[bookarrow] (e) -- (f);
\node[booknote] at (5.6,-.9) {不同身份各自进入判断；账户排队决定资金条件，输入相同不合并身份};
\end{tikzpicture}
\caption{同参不同身份分别判断，初始资金决定两成或一成一拒}
\label{fig:ch06-distinct-identities}
\end{figure}


原结果一致的证明从发布时开始。首次提交的 \(R(S)\) 由该事务实际 P、O、D 和成交值构造，编号、金额与当次余额通过业务关联不变量核对，因而它描述真实终态。以后所有同参恢复只读取这份结果，终态保存内容不可变，结果版本解释保持，所以每次业务投影都等于同一 \(R(S)\)。这同时保证稳定与真实关联，避免把一个永远不变的错误 D9 当作一致性已经足够。

本次 \code{Date} 或 \code{request\_id} 的变化不进入这项比较。\code{Date} 表示响应产生的时间，\code{request\_id} 标识当前通信；原结果中的 \code{balance\_after\_minor} 则进入比较，因为它描述历史扣款。若 S2 先后完成使当前余额变为 4000，恢复 S1 仍需返回原当次 7000。把当前余额当成原字段，记录数量与资金守恒可能全部正确，结果条件仍会独立失败。

输入冲突还保护结果的解释。同一 \(S\) 已绑定 Q7，后来 Q8 到达，应返回 409 并保留旧结果；随后 Q7 正确恢复仍能得到 \(R(S)\)。若把冲突写成 \(S\) 的新 \code{failed}，会破坏不可变终态；若直接把原成功当成 Q8 的成功，又破坏输入绑定。先用旧规范化规则比较，再选择返回原结果或当前冲突，才保持身份、输入与结果指向同一项操作。
\FloatBarrier
